% ============================================================================
% Section 9.7 — Simulating Compound Events
% CCSS 7.SP.C.8c
% ============================================================================
\section{Simulating Compound Events}

\begin{stepGoal}
\begin{itemize}
    \item Design a simulation using random numbers or physical tools
    \item Use simulation results to estimate probabilities of compound events
\end{itemize}
\end{stepGoal}

\begin{stepVocabBox}
    \vocabItem{Simulation}{Using a simple model (coins, dice, random digits) to \textbf{imitate} a real-world situation.}
    \vocabItem{Random Digits}{Numbers $0$--$9$ where each digit is equally likely to appear.}
\end{stepVocabBox}

\begin{stepsCard}[How to Design and Run a Simulation]
    \stepItem{Identify the \textbf{event} and its real-world probability.}
    \stepItem{Choose a \textbf{model} (coin, die, spinner, or random digits) that matches the probability.}
    \stepItem{Define what \textbf{one trial} looks like and run \textbf{many trials} (at least $20$--$50$).}
    \stepItem{Record results, calculate \textbf{experimental probability}, and compare to what you expect.}
\end{stepsCard}

% --- Worked Example 1 ---
\begin{stepExample}{About $50\%$ of customers order fries. Simulate $10$ customers by flipping a coin. Here are the results: $H, T, H, H, T, T, H, T, H, H$. Estimate the probability that a customer orders fries.}
    \stepShow{1} Event: customer orders fries. Probability: $50\%$.

    \stepShow{2} Model: coin flip ($H$ = fries, $T$ = no fries) since $P(H) = 0.5$ matches.

    \stepShow{3} One trial = one flip per customer. Run $10$ trials.\par
    Results: $H, T, H, H, T, T, H, T, H, H$.

    \stepShow{4} Fries (heads): $6$ out of $10$.\par
    $P(\text{fries}) = \dfrac{6}{10} = 0.6$. Expected: $0.5$. Close --- more trials would get closer.
    \stepResult{Experimental $P \approx 0.6$; expected $0.5$.}
\end{stepExample}

% --- Worked Example 2 ---
\begin{stepExample}{A basketball player makes $70\%$ of free throws. Use random digits ($0$--$9$) to simulate $5$ sets of $2$ free throws. Let digits $0$--$6$ = make, $7$--$9$ = miss. Random digits: $3\;8 \mid 1\;5 \mid 9\;2 \mid 0\;7 \mid 4\;6$. How many times does the player make both shots?}
    \stepShow{1} Event: make a free throw. Probability: $70\%$.

    \stepShow{2} Model: digits $0$--$6$ = make ($7$ digits out of $10 = 70\%$); digits $7$--$9$ = miss.

    \stepShow{3} One trial = $2$ consecutive digits. Five trials:\par
    Trial~1: $3, 8$ $\to$ make, miss\par
    Trial~2: $1, 5$ $\to$ make, make \checkmark\par
    Trial~3: $9, 2$ $\to$ miss, make\par
    Trial~4: $0, 7$ $\to$ make, miss\par
    Trial~5: $4, 6$ $\to$ make, make \checkmark

    \stepShow{4} Made both: $2$ out of $5$ trials.\par
    $P(\text{both make}) = \dfrac{2}{5} = 0.4$. Theoretical: $0.7 \times 0.7 = 0.49$. Fairly close for just $5$ trials.
    \stepResult{Experimental: $0.4$; Theoretical: $0.49$}
\end{stepExample}

\watchOut{Make sure your model matches the real probability. For $70\%$, use $7$ out of $10$ digits --- not $7$ out of $9$!}

% ============================================================================
% PRACTICE
% ============================================================================
\newpage
\begin{stepPractice}{Simulating Compound Events Practice}
    \resetProblems

    \stepPracticeHeader[step1Color]{Choose a Model}
    \prob You want to simulate a $30\%$ chance of rain. Using digits $0$--$9$, which digits represent rain? \answerBlank[3cm]
    \answer{$0, 1, 2$ (any $3$ out of $10$ digits)}

    \stepPracticeHeader[step2Color]{Run the Simulation}
    \prob Use these random digits for $10$ trials of the $30\%$ rain model: $5, 0, 8, 2, 7, 1, 4, 9, 3, 6$. How many rain events? \answerBlank[2cm]
    \answer{$4$ (digits $0, 2, 1, 3$)}

    \stepPracticeHeader[step3Color]{Calculate and Compare}
    \prob What is the experimental probability of rain from the $10$ trials above? How does it compare to $30\%$? \answerBlank[3cm]
    \answer{$\frac{4}{10} = 40\%$; slightly above $30\%$ but reasonable for $10$ trials}

    \prob A spinner is $50\%$ blue and $50\%$ red. You spin it twice. Simulate $20$ trials using coin flips. If you got ``both red'' $6$ times, what is the experimental probability? What is the theoretical? \answerBlank[3cm]
    \answer{Experimental: $\frac{6}{20} = 0.3$; Theoretical: $0.5 \times 0.5 = 0.25$}

    \wordProblem{A store gives a scratch card to each customer. Each card has a $20\%$ chance of winning a prize. Using digits $0$--$9$, design a simulation to estimate the probability that $2$ customers in a row both win. Describe your model and define one trial.}{~}
    \answer{Let digits $0, 1$ = win ($2$ out of $10 = 20\%$). One trial = two consecutive random digits. ``Both win'' if both digits are $0$ or $1$. Repeat many trials and divide successes by total trials. Theoretical: $0.2 \times 0.2 = 0.04$.}
\end{stepPractice}
